% A 3-5 dislocation between two parallel turns (Propositions 3 and 6).
% French copy of dislocation.tex: only the labels differ.
% \input inside a figure; needs \usepackage{tikz} and
% \usetikzlibrary{arrows.meta,decorations.pathreplacing}.
%
% Each panel is a band of a T-mesh on a cylinder, cut open along a trace of
% the vertex c_5 of valence five: the right side is glued to the left side, the
% lower part (below c_5) at the same height, the upper part two units higher.
% Every node is at a whole position and the faint grid is the quad mesh at
% these lengths. c_5 is drawn three times (its angle 5 x 90 degrees is split
% between the three copies: 90 at the bottom left, 180 at the right, 180 at
% the left), c_3 once (270 degrees). The empty corner next to c_3 is not part of
% the surface: its two thick sides are glued to each other by a quarter-turn
% about c_3. The turn gamma_1 (below the pair) and gamma_2 (above it) bound an
% annulus that holds exactly c_3 and c_5 (Proposition 2(b)).
%
% (a) Straight chain (Proposition 3(ii)). c_3 = (2,3), c_5 = (0,3) = (10,3) =
% (2,5); the gluings are (10,y) ~ (0,y) for y < 3, (10,y) ~ (2,y+2) for y > 3,
% and the quarter-turn about c_3 taking [c_3, (0,3)] onto [c_3, (2,5)]. c_3 and
% c_5 are joined along the trace y = 3 by the chain of arcs b = [c_3, (1,3)]
% and c = [(1,3), c_5]; a trace from below (its vertex is out of the picture)
% stops at (1,3). delta = l_b + l_c = 1 + 1 = 2.
%   gamma_1 at o = 1.5: crosses x = 1, 2, 10; comes back at 1.5: s = 0,
%   length 10 (a ring).
%   gamma_2 at o = 5.5 on the side x = 2: crosses x = 10 at height 5.5, which
%   is height 7.5 on the side x = 2: s = 2, length 10 - 2 = 8 (a helix).
%   |s_2 - s_1| = 2 = delta (Proposition 3(ii)); the lengths differ by 2 too:
%   the Burgers vector of the pair is (2, 2), of norm sqrt(2) |d| with d = (2, 0)
%   from c_3 to c_5 (Proposition 3(i)); s_2 - s_1 = 2 = d_1 + d_2 mod 2 (iii).
%   The separatrix y = 5 of c_5 goes once round and stops on the trace x = 10
%   two units above c_5 (Proposition 8 gives s_2 = 2 again).
%   Both turns can be rings only if delta = 0: c_3 and c_5 merge.
%   Parity on a sphere (Proposition 5): 10 = 0 and 8 = 2 mod 2.
%
% (b) Staircase (Proposition 6). c_3 = (1,4), c_5 = (0,3) = (10,3) = (0,5);
% gluings (10,y) ~ (0,y) for y < 3, (10,y) ~ (0,y+2) for y > 3, and the
% quarter-turn about c_3 taking [c_3, (1,3)] onto [c_3, (0,4)] and
% [(1,3), (0,3)] onto [(0,4), (0,5)]. c_3 and c_5 are on no common trace; the
% staircase from c_3 to c_5 is b = [c_3, (1,3)], down, then c = [(1,3), c_5],
% left: d = p_5 - p_3 = (-1, -1), and with epsilon = (-1, -1) each step has
% epsilon . v = 1 (b: v = (0,-1); c: v = (-1,0)).
%   gamma_1 at o = 1.5: s = 0, length 10. gamma_2 at o = 5.5: crosses x = 10
%   at height 5.5 = height 7.5 on the side x = 0: s = 2, length 10.
%   s_2 - s_1 = 2 = l_b + l_c = epsilon . d (Proposition 6); here the lengths
%   are equal, the Burgers vector of the pair is (0, 2). Both are rings only
%   if l_b = l_c = 0. Parity: 10 = 0 = 2 mod 2.
% check.py builds both bands and checks these numbers by walking the strips
% of the quad mesh: valences 5 and 3 at the copies listed above, every other
% vertex of valence 4, and the turns above.
\begin{tikzpicture}[x=0.5cm, y=0.5cm, font=\footnotesize,
    side/.style={draw=black, line width=0.5pt},
    quadgrid/.style={draw=black!18, line width=0.25pt, step={(1,1)}},
    surface/.style={fill=black!3},
    node dot/.style={circle, fill=black, inner sep=0pt, minimum size=2.6pt},
    arc mark/.style={line width=1.8pt, draw=black},
    turn/.style={line width=1pt, -{Stealth[length=5pt]}},
    brace/.style={decorate, decoration={brace, amplitude=3pt}, line width=0.5pt},
    val3/.pic={\fill[valthree, draw=black, line width=0.3pt]
      (90:3.4pt) -- (210:3.4pt) -- (330:3.4pt) -- cycle;},
    val5/.pic={\fill[valfive, draw=black, line width=0.3pt]
      (90:3.4pt) -- (162:3.4pt) -- (234:3.4pt) -- (306:3.4pt) -- (18:3.4pt) -- cycle;}]
  \definecolor{loopspiral}{HTML}{D55E00}
  \definecolor{loopclosed}{HTML}{0B3C8C}
  \definecolor{valthree}{HTML}{E8890C}
  \definecolor{valfive}{HTML}{7C3AAD}

  % ---------------------------------------------------------------- (a)
  \begin{scope}
    % The surface: the lower band [0,10]x[0,3] and the upper part [2,10]x[3,8].
    \fill[surface] (0,-0.6) rectangle (10,3) (2,3) rectangle (10,8.6);
    \draw[quadgrid] (0,-0.6) grid (10,3);
    \draw[quadgrid] (2,3) grid (10,8.6);
    % Traces: the cut (a trace of c_5, both copies), the trace from c_3 round
    % to c_5, c_3 downwards, a trace from below, the winding separatrix.
    \draw[side] (0,-0.6) -- (0,3) (10,-0.6) -- (10,8.6) (2,5) -- (2,8.6);
    \draw[side] (2,3) -- (10,3) (2,-0.6) -- (2,3) (1,-0.6) -- (1,3) (2,5) -- (10,5);
    % The chain from c_3 to c_5, and its glued copy.
    \draw[arc mark] (0,3) -- (2,3) (2,3) -- (2,5);
    \node[below] at (1.5,3) {$b$};
    \node[below] at (0.5,3) {$c$};
    \foreach \p in {(1,3), (2,4), (10,5), (2,7)} \node[node dot] at \p {};
    % The turns.
    \draw[turn, loopclosed] (0,1.5) -- (10.1,1.5);
    \draw[loopclosed, line width=0.5pt, fill=white] (0,1.5) circle (1.8pt);
    \node[above, loopclosed] at (6,1.5) {$\gamma_1$};
    \draw[turn, loopspiral] (2,5.5) -- (10.1,5.5);
    \draw[loopspiral, line width=0.5pt, fill=white] (2,5.5) circle (1.8pt);
    \draw[loopspiral, dashed, line width=0.8pt, -{Stealth[length=4pt]}] (2,7.5) -- (4,7.5);
    \draw[loopspiral, line width=0.5pt, fill=white] (2,7.5) circle (1.8pt);
    \node[above, loopspiral] at (6,5.5) {$\gamma_2$};
    \draw[brace] (1.7,5.5) -- (1.7,7.5) node[midway, left=3pt] {$2$};
    % The singular vertices.
    \pic at (2,3) {val3};
    \foreach \p in {(0,3), (10,3), (2,5)} \pic at \p {val5};
    \node[above right=-1pt] at (2,3) {$c_3$};
    \node[left=1pt] at (0,3) {$c_5$};
    % Right: shifts and lengths.
    \node[right, align=left, loopspiral] at (10.3,5.5) {$s = 2$\\ 8 quads};
    \node[right, align=left, loopclosed] at (10.3,1.5) {$s = 0$\\ 10 quads};
    \node[anchor=north, align=center] at (6,-1.1)
      {$s_{\gamma_2} - s_{\gamma_1} = \ell_b + \ell_c = \delta = 2$};
    \node[anchor=north] at (6,-2.4) {\small (a) chaîne droite};
  \end{scope}

  % ---------------------------------------------------------------- (b)
  \begin{scope}[shift={(16.8,0)}]
    % The surface: [0,10]x[0,3], [1,10]x[3,4], [0,10]x[4,8].
    \fill[surface] (0,-0.6) rectangle (10,3) (1,3) rectangle (10,4) (0,4) rectangle (10,8.6);
    \draw[quadgrid] (0,-0.6) grid (10,3);
    \draw[quadgrid] (1,3) grid (10,4);
    \draw[quadgrid] (0,4) grid (10,8.6);
    % Traces: the cut (a trace of c_5), the loop of c_5 at y = 3, the three
    % separatrices of c_3, the winding separatrix of c_5.
    \draw[side] (0,-0.6) -- (0,3) (10,-0.6) -- (10,8.6) (0,5) -- (0,8.6);
    \draw[side] (0,3) -- (10,3) (1,4) -- (10,4) (1,4) -- (1,5) (0,5) -- (10,5);
    % The staircase from c_3 to c_5, and its glued copy.
    \draw[arc mark] (1,4) -- (1,3) -- (0,3) (1,4) -- (0,4) -- (0,5);
    \node[right] at (1,3.5) {$b$};
    \node[below] at (0.5,3) {$c$};
    \foreach \p in {(1,3), (0,4), (10,4), (0,6), (1,5), (10,5), (0,7)}
      \node[node dot] at \p {};
    % The turns.
    \draw[turn, loopclosed] (0,1.5) -- (10.1,1.5);
    \draw[loopclosed, line width=0.5pt, fill=white] (0,1.5) circle (1.8pt);
    \node[above, loopclosed] at (6,1.5) {$\gamma_1$};
    \draw[turn, loopspiral] (0,5.5) -- (10.1,5.5);
    \draw[loopspiral, line width=0.5pt, fill=white] (0,5.5) circle (1.8pt);
    \draw[loopspiral, dashed, line width=0.8pt, -{Stealth[length=4pt]}] (0,7.5) -- (2,7.5);
    \draw[loopspiral, line width=0.5pt, fill=white] (0,7.5) circle (1.8pt);
    \node[above, loopspiral] at (6,5.5) {$\gamma_2$};
    \draw[brace] (-0.3,5.5) -- (-0.3,7.5) node[midway, left=3pt] {$2$};
    % The singular vertices.
    \pic at (1,4) {val3};
    \foreach \p in {(0,3), (10,3), (0,5)} \pic at \p {val5};
    \node[above right=-1pt] at (1,4) {$c_3$};
    \node[left=1pt] at (0,3) {$c_5$};
    % Right: shifts and lengths.
    \node[right, align=left, loopspiral] at (10.3,5.5) {$s = 2$\\ 10 quads};
    \node[right, align=left, loopclosed] at (10.3,1.5) {$s = 0$\\ 10 quads};
    \node[anchor=north, align=center] at (6,-1.1)
      {$s_{\gamma_2} - s_{\gamma_1} = \ell_b + \ell_c = 2$};
    \node[anchor=north] at (6,-2.4) {\small (b) escalier};
  \end{scope}
\end{tikzpicture}
